\section{Biologia}

\begin{itemize}
  \item \textbf{evoluzione}: le popolazioni di organismi si sono evolute nel tempo a partire da forme di vita primordiali;
  \item \textbf{trasmissione dell'informazione}: l'evoluzione dipende dalla trasmissione dell'informazione genetica da una generazione all'altra;
  \item \textbf{trasferimento dell'energia}: tutti i processi vitali richiedono un continuo ingresso di energia.
\end{itemize}

% ---------------------------------------------------------------------------
\subsection{I livelli di organizzazione della vita}

Gerarchia dei livelli, dal più semplice al più complesso:
\[
\text{atomo} \rightarrow \text{molecola} \rightarrow \text{macromolecola} \rightarrow \text{procariote} \rightarrow \text{organello} \rightarrow \text{cellula eucariote} \rightarrow \text{tessuto}
\]
\[
\rightarrow \text{organo} \rightarrow \text{apparato} \rightarrow \text{organismo} \rightarrow \text{popolazione} \rightarrow \text{comunità} \rightarrow \text{ecosistema} \rightarrow \text{biosfera}.
\]

\rubrica{I viventi: livelli di organizzazione}
\begin{center}
\begin{forest} classalbero
[I viventi
  [Livello molecolare [Virus][Viroidi][Prioni]]
  [Livello cellulare
    [Procarioti [Archea][Batteri]]
    [Eucarioti [Protozoi][Lieviti][Cellule animali][Cellule vegetali]]
  ]
]
\end{forest}
\end{center}

% ---------------------------------------------------------------------------
\subsection{La classificazione dei viventi}

\textbf{Tre domini} (dall'analisi delle sequenze di rRNA di Carl Woese; tutti derivati da un progenitore comune):
\begin{multicols}{3}
  \begin{itemize}
    \item \textbf{Bacteria} (Batteri);
    \item \textbf{Archaea} (Archei);
    \item \textbf{Eukarya} (Eucarioti).
\end{itemize}
\end{multicols}

\textbf{Sei regni}: Bacteria, Archaea, Protista, Plantae, Animalia, Fungi.

\rubrica{Gerarchia di classificazione linneana}
Dal raggruppamento più ampio al più specifico:
\[
\text{dominio} \rightarrow \text{regno} \rightarrow \text{phylum} \rightarrow \text{classe} \rightarrow \text{ordine} \rightarrow \text{famiglia} \rightarrow \text{genere} \rightarrow \text{specie}.
\]
Esempi:
\begin{itemize}
  \item Eukarya $\rightarrow$ Animalia $\rightarrow$ Chordata $\rightarrow$ Mammalia $\rightarrow$ Primates $\rightarrow$ Pongidae $\rightarrow$ \textit{Pan} $\rightarrow$ \textit{Pan troglodytes};
  \item Eukarya $\rightarrow$ Animalia $\rightarrow$ Artropodi $\rightarrow$ Insetti $\rightarrow$ Blattoidei $\rightarrow$ Blattellidi $\rightarrow$ \textit{Blattella} $\rightarrow$ \textit{Blattella germanica}.
\end{itemize}

\textbf{Specie}: insieme di popolazioni interfertili fra loro e isolate riproduttivamente dalle altre.

% ---------------------------------------------------------------------------
\subsection{Le teorie evoluzionistiche}

\rubrica{Linneo (fine 1700)}
Fondatore della moderna classificazione dei viventi; definiva le specie come ``fisse'', cioè invariabili nel tempo e incapaci di modificarsi. All'inizio del XIX secolo la teoria fu messa in discussione: negli strati rocciosi più antichi mancavano tracce fossili delle specie allora viventi, mentre se ne rinvenivano altre appartenenti a organismi estinti.

\rubrica{Lamarck (1809, \textit{Philosophie Zoologique})}
Prima teoria evoluzionista: ipotizzò la modificabilità delle specie sotto l'influenza delle condizioni ambientali.
\begin{itemize}
  \item principio: è il bisogno di una funzione che crea l'organo; l'uso o il non uso di determinati organi porta con il tempo a un loro potenziamento o a una loro atrofia;
  \item errore di fondo: l'ereditabilità dei caratteri acquisiti.
\end{itemize}

\rubrica{Darwin (1859, \textit{L'origine delle specie per mezzo della selezione naturale})}
Nuova teoria evolutiva basata su tre concetti fondamentali: \textbf{variabilità}, \textbf{selezione naturale}, \textbf{ereditarietà}.
\begin{itemize}
  \item viaggio del \textit{Beagle} (1831--1836): circumnavigazione del globo;
  \item albero della vita: la storia della vita immaginata come un albero; i punti di ramificazione rappresentano l'origine di nuove famiglie, i rami che non raggiungono il margine superiore rappresentano gruppi estinti.
\end{itemize}
Osservazioni e deduzioni:
\begin{enumerate}
  \item la maggior parte degli organismi produce più di uno o due figli; le popolazioni non aumentano indefinitamente di dimensioni; cibo e risorse sono limitati $\rightarrow$ gli individui competono per risorse limitate;
  \item gli individui mostrano variabilità in molte caratteristiche; molte variazioni hanno basi genetiche ereditate $\rightarrow$ le caratteristiche ereditarie permettono ad alcuni individui di sopravvivere più a lungo e riprodursi più di altri;
  \item conclusione: le caratteristiche di una popolazione cambieranno nel corso delle generazioni perché quelle vantaggiose ed ereditabili diventeranno più comuni.
\end{enumerate}

% ---------------------------------------------------------------------------
\subsection{La teoria cellulare}

Tre principi:
\begin{itemize}
  \item le cellule sono le unità fondamentali della vita;
  \item tutti gli organismi sono composti da cellule;
  \item tutte le cellule provengono da cellule preesistenti.
\end{itemize}

\rubrica{Scoperta delle cellule}
\begin{itemize}
  \item \textbf{Robert Hooke} (1665): con un microscopio composto osserva, in una sezione sottile di sughero, un reticolo di ``celle'';
  \item \textbf{Anton van Leeuwenhoek} (1673): con un microscopio a lente singola (ingrandimento $\sim$270 volte, risoluzione $\sim$1,35\,$\mu$m) osserva batteri e altri microrganismi.
\end{itemize}

% ---------------------------------------------------------------------------
\subsection{Ordini di grandezza e microscopia}

\rubrica{Limite di risoluzione}
Distanza minima al di sotto della quale non siamo più in grado di vedere due punti distinti. Definisce l'ambito d'impiego dei tre strumenti:
\begin{itemize}
  \item \textbf{occhio umano}: strutture dal millimetro/centimetro in su;
  \item \textbf{microscopio ottico}: scala del micrometro;
  \item \textbf{microscopio elettronico}: scala del nanometro.
\end{itemize}

Dimensioni di riferimento:
\begin{table}[htbp]
  \centering \small \renewcommand{\arraystretch}{1.3}
  \begin{tabularx}{\textwidth}{|X|l|l|}
    \hline
    \textbf{Struttura} & \textbf{Dimensione} & \textbf{Strumento} \\
    \hline
    Uovo di gallina (tuorlo) & 3 cm & occhio nudo \\
    \hline
    Uova di rana o pesce & 1 mm & occhio nudo \\
    \hline
    Ovulo umano & 100\,$\mu$m & microscopio ottico \\
    \hline
    Tipica cellula vegetale & 10--100\,$\mu$m & microscopio ottico \\
    \hline
    Tipica cellula animale & 5--30\,$\mu$m & microscopio ottico \\
    \hline
    Cloroplasto & 2--10\,$\mu$m & microscopio ottico \\
    \hline
    Mitocondrio & 1--5\,$\mu$m & microscopio ottico \\
    \hline
    \textit{Escherichia coli} & 1\,$\mu$m & microscopio ottico \\
    \hline
    Grandi virus (HIV, influenza) & 100 nm & microscopio elettronico \\
    \hline
    Ribosoma & 25 nm & microscopio elettronico \\
    \hline
    Membrana cellulare (spessore) & 7--10 nm & microscopio elettronico \\
    \hline
    DNA a doppia elica (diametro) & 2 nm & microscopio elettronico \\
    \hline
    Atomo d'idrogeno & 0,1 nm & microscopio elettronico \\
    \hline
  \end{tabularx}
\end{table}

\rubrica{Microscopio ottico}
Componenti (dal basso): sorgente luminosa $\rightarrow$ condensatore $\rightarrow$ campione $\rightarrow$ lente dell'obiettivo $\rightarrow$ lente dell'oculare. Permette di vedere cellule colorate o vive, ma a risoluzione relativamente bassa.

Tecniche di microscopia ottica:
\begin{itemize}
  \item \textbf{campo luminoso}: la luce passa direttamente attraverso il campione; molte strutture hanno contrasto troppo basso $\rightarrow$ si usa un colorante, che però di norma fissa e uccide le cellule;
  \item \textbf{campo scuro}: il campione è illuminato con una inclinazione specifica e solo la luce da esso diffusa raggiunge la lente $\rightarrow$ immagine luminosa su sfondo nero;
  \item \textbf{contrasto di fase}: le differenze di rifrazione dovute alla densità del campione sono rese come differenze di contrasto $\rightarrow$ evidenzia strutture altrimenti invisibili; le cellule vive possono essere fotografate o filmate;
  \item \textbf{Nomarski} (contrasto a interferenza differenziale): lenti speciali amplificano le differenze di densità $\rightarrow$ aspetto tridimensionale;
  \item \textbf{fluorescenza}: strutture o molecole sono marcate con coloranti fluorescenti che emettono luce sotto illuminazione ultravioletta $\rightarrow$ localizzazione delle molecole marcate;
  \item \textbf{confocale a scansione laser}: un laser scandisce il campione marcato e il computer focalizza la luce su un singolo piano $\rightarrow$ immagine tridimensionale più nitida delle altre tecniche ottiche.
\end{itemize}

\rubrica{Microscopio elettronico}
\begin{itemize}
  \item \textbf{a trasmissione (MET)}: il fascio elettronico attraversa il campione $\rightarrow$ immagine ad altissima risoluzione, fortemente ingrandita (solo una parte di una sezione sottile);
  \item \textbf{a scansione (MES)}: sfrutta gli elettroni secondari $\rightarrow$ chiara visione delle caratteristiche della superficie della cellula.
\end{itemize}
A parità di ingrandimento (4500 volte), il microscopio elettronico mostra molti più dettagli del microscopio ottico (miofibrille, mitocondri, membrane delle singole strutture): il microscopio ottico non può fornire informazioni migliori.
